\section{Background}

\subsection{Safety Cases and GSN}
A safety case is a structured argument that a system is acceptably
safe. The Goal Structuring Notation (GSN)~\cite{kelly1998,kelly2004}
is the de facto standard for expressing such arguments. A GSN
argument is a directed acyclic graph of goals (claims), strategies
(reasoning steps), and solutions (evidence). For UAVs, EASA
SC-VTOL~\cite{easa_sc_vtol} and DO-178C~\cite{do178c} define the
regulatory context.

The limitation of GSN in practice is that solutions are static
artefacts: a test report, a simulation log, a hardware-in-the-loop
trace. When the system changes, the solutions become stale, but the
GSN argument is rarely updated. This disconnect is the problem SCC
addresses.

\subsection{Runtime Verification}
Runtime verification~\cite{bartocci2018} monitors a running system
against a formal specification and raises an alarm when the
specification is violated. It is complementary to static safety
cases: the safety case argues \emph{why} the system should be safe;
runtime verification checks \emph{whether} it is.

SCC bridges the two: it compiles a safety requirement into both a
runtime monitor (a Rust test) and a GSN node, ensuring that the
argument and the monitor are generated from the same source.

\subsection{Hash Chains for Audit}
A hash chain~\cite{crosby2009,merkle1988} is a sequence of blocks
where each block contains the hash of the previous one. Any
modification to an earlier block invalidates all subsequent hashes,
making tampering detectable. SCC uses a hash chain to record each
compiled requirement, providing a tamper-evident audit trail that
a regulator can verify independently of the operator.

\section{SCC Design}

\subsection{Requirements DSL}
SCC accepts a small domain-specific language for UAV safety
requirements. Each line has the form:
\begin{verbatim}
<name>: <value(s)> <severity>
\end{verbatim}
where \texttt{severity} is one of \texttt{HIGH}, \texttt{MEDIUM},
\texttt{LOW}. Three requirement types are supported:
\begin{itemize}
    \item \texttt{speed\_max}: a scalar upper bound on speed.
    \item \texttt{altitude\_range}: a lower and upper bound on altitude.
    \item \texttt{obstacle\_clearance}: a minimum distance to obstacles.
\end{itemize}
Example:
\begin{verbatim}
speed_max: 10.0 HIGH
altitude_range: 5.0 120.0 HIGH
obstacle_clearance: 2.0 MEDIUM
\end{verbatim}

\subsection{Compilation Pipeline}
The pipeline has four stages:
\begin{enumerate}
    \item \textbf{Parse}: the DSL is parsed into a list of
          \texttt{Requirement} structs. Malformed lines are rejected
          with a line number and reason.
    \item \textbf{Generate tests}: each requirement becomes a Rust
          test stub that asserts the invariant.
    \item \textbf{Generate GSN}: each requirement becomes a GSN node
          with severity, evidence reference, and chain block index.
    \item \textbf{Compile chain}: each requirement becomes a block in
          a SHA-256 hash chain, binding the requirement to the
          mission context.
\end{enumerate}
The output is three artefacts per requirement, all derived from the
same input. There is no manual synchronisation.

\subsection{Why Rust}
Rust was chosen for three reasons: (i) the flight control stack
(\texttt{rnfc-core}) is already in Rust, so SCC integrates
natively; (ii) the \texttt{no\_std} subset allows SCC to run on the
companion computer without a full OS; (iii) the type system makes
the parser and the code generator easy to test exhaustively.

\section{Implementation}

\subsection{Crate Structure}
SCC is implemented as the crate \texttt{rnfc-scc} (160 lines of
Rust). The public API has four functions:
\begin{itemize}
    \item \texttt{parse\_requirements(input: \&str) -> Result<Vec<Requirement>, String>}
    \item \texttt{generate\_test(req: \&Requirement) -> String}
    \item \texttt{generate\_gsn(req: \&Requirement, block: usize) -> String}
    \item \texttt{compile\_requirements(reqs: \&[Requirement]) -> Result<CertChain, String>}
\end{itemize}
The crate depends on \texttt{rnfc-certchain} (SHA-256 hash chain)
and uses \texttt{sha2} for the hash function.

\subsection{Integration with CertChain}
Each requirement is hashed together with its values and severity,
and the hash is appended to a \texttt{CertChain}. The chain is
deterministic: the same requirements always produce the same chain
root. This allows a regulator to verify that a drone was certified
against a specific set of requirements, and that the requirements
have not changed since certification.

\section{Evaluation}

\subsection{Case Study: 3-Requirement UAV Envelope}
We evaluate SCC on the envelope: \texttt{speed\_max = 10.0~m/s},
\texttt{altitude\_range = [5.0, 120.0]~m}, \texttt{obstacle\_clearance
= 2.0~m}. All three are \texttt{HIGH} or \texttt{MEDIUM} severity.

\subsection{Results}
The compiler produces:
\begin{itemize}
    \item 3 Rust test stubs (one per requirement), each with the
          assertion and the severity as a comment.
    \item 3 GSN nodes, each with the requirement, severity,
          evidence reference (\texttt{test\_<name>\_invariant}), and
          chain block index.
    \item A 3-block \texttt{CertChain} with a deterministic root hash.
\end{itemize}
A single-line change (\texttt{speed\_max: 10.0} to \texttt{8.0})
regenerates all three artefacts and produces a different chain root.
The change is therefore visible in the audit trail.

\subsection{Limitations}
SCC supports only three requirement types. It does not support
quantifiers (\emph{for all} trajectories), temporal operators
(\emph{always}, \emph{eventually}), or composition of safety cases.
The physics used to evaluate the requirements is synthetic (1D/2D),
not 6-DOF. These are open problems.

\section{Discussion}

\subsection{Relation to DO-326A}
DO-326A~\cite{do326a} requires a security risk assessment and a
plan for security aspects of certification (PSecAC). SCC provides
evidence for the \emph{safety} side of the assessment: each
requirement is compiled into a test and a chain block, and the
chain is anchored externally (via \texttt{rnfc-anchor}). The
security requirements (SEC-01 to SEC-13) are tracked separately in
the project's \texttt{SECURITY\_REQUIREMENTS.md}.

\subsection{Open Problems}
Three problems remain open: (i) a formal semantics for the DSL, so
that the generated tests can be proven equivalent to the
requirements; (ii) composition, so that multiple safety cases can
be combined; (iii) integration with runtime monitors, so that the
compiled tests are not only stubs but active monitors on the drone.

\section{Conclusion}
We presented SCC, a tool that compiles safety requirements into
executable evidence: Rust tests, GSN nodes, and hash chain blocks.
The key property is single source of truth: changing one line in
the requirements regenerates all artefacts and re-validates the
chain. SCC is implemented in Rust, tested with 10 unit tests, and
evaluated on a 3-requirement UAV envelope. Future work includes
formal semantics, composition, and runtime monitoring.
